% !PRINTED 498
\weberpage{498}
%<<CONT
and one can therefore, if besides the four figures 0, 1, 2, 3 a fifth, 4, is present, compose all three-term cycles from pairs of transpositions of the form $(0, 1)(3, 4)$.

To this we append the proofs of two further important theorems relating to transitive permutation groups:

\gesperrt{9. If a transitive permutation group of $m$ figures contains a single transposition, then it is either the symmetric group or it is imprimitive.}

Let the given group be $P$, and let the $m$ figures that are permuted by it be $0, 1, 2 \dots m - 1$, and let the transposition $(0, 1)$ occur in it. Suppose that it contains altogether the transpositions
\[
(0, 1), (0, 2), (0, 3), \dots (0, \mu - 1),
\]
but no other transposition with the figure 0. Now if $\mu = m$, then by 2. $P$ is the symmetric group. But if $\mu < m$, then by the same theorem $P$ contains all permutations of the figures
\[
M = 0, 1, 2 \dots \mu - 1.
\]
From this it follows that $P$ can contain no transposition of one of the figures $0, 1, 2, 3 \dots \mu - 1$ with another figure, say with $\mu$. For if $(1, \mu)$ occurs in $P$, then $(0, \mu) = (0, 1)(1, \mu)(0, 1)$ also occurs in it, contrary to the hypothesis.

Now if $P$ is transitive, there is in it a permutation $\pi$ by which the figure 0 passes into a figure not contained in the system $M$, say into $\mu$; and then by $\pi$ none of the figures of $M$ can pass into a figure of $M$; for if, e.\,g., 1 passes into $r$, then
\[
\pi^{-1}(0, 1)\pi = (\mu, r)
\]
and $(\mu, r)$, if $r$ belongs to $M$, cannot, as we have seen, occur in $P$. By every permutation of $P$, therefore, the figures of the system $M$ are either permuted only among themselves, or they are carried over into the figures of a system $M'$ entirely distinct from it, and in $P$ all permutations of the figures $M'$ also occur.

If with $M$ and $M'$ not all the figures are yet exhausted, there is a permutation in $P$ by which $M$ is carried over into a third system $M''$, which again can have no figure in common with $M$ or with $M'$.
